\selectlanguage{ngerman}%
\frenchspacing%
\begin{ledgroupsized}[r]{120mm}%
\footnotesize%
\pstart%
\noindent\textbf{Überlieferung:}
\pend%
\end{ledgroupsized}%
\begin{ledgroupsized}[r]{114mm}%
\footnotesize%
\pstart%
\parindent -6mm%
\makebox[6mm][l]{\textit{L}}%
Notiz:
LH XXXV 11,\;14 Bl.\;37.
Ein Zettel (6\;x\;5,5 cm);
Ränder beschnitten.
Eine voll beschriebene Seite auf Bl.\;\href{\Scan{LH_35_11_14_037v}}{37\:v\textsuperscript{o}}.
Bl.\;\href{\Scan{LH_35_11_14_037r}}{37\:r\textsuperscript{o}} überliefert lediglich folgende Textfragmente von Leibnizens Hand:
\textls{\textit{Ars combinator}}\textlangle\textls{\textit{ia}}\textrangle\ \textlangle...\textrangle\ \lbrack/\rbrack\ \textlangle \textit{Libell}\textrangle\textit{us ab autore admodum} \textlangle\textit{juvene}\textrangle\ \textlangle...\textrangle\ \lbrack/\rbrack\ \textlangle--\textrangle\ \sout{\textit{ob cogitationem}} \textit{ob nova} \textlangle...\textrangle\ \lbrack/\rbrack\ \textlangle--\textrangle\ \textit{artis inveniendi semina} \textlangle...\textrangle\ \lbrack/\rbrack\ \textlangle--\textrangle\textit{dicio usu atque annis elima} \textlangle...\textrangle\
Hierbei handelt es sich wohl um einen Entwurf seiner im Februar 1691 anonym veröffentlichten Besprechung der unauthorisierten Neuausgabe seiner Jugendschrift \textit{Dissertatio de arte combinatoria}
(vgl. \textit{AE} 1691, S.\;63\:f.; \textit{LSB} VI,\;2, S.\;549.9--550.5;\cite{01023}\cite{01486}
siehe zudem \textsc{G.\:W.\;Leibniz}, \textit{Ars combinatoria}, Frankfurt a.\:M.\;1690).\cite{01485}
\pend%
\end{ledgroupsized}%
\vspace{5mm}%
\begin{ledgroup}%
\footnotesize%
\pstart%
\noindent%
\textbf{Datierungsgründe:}
Die Datierung der Notiz N.\;\ref{58952} beruht auf Eigenschaften des Textträgers.
Siehe die ausführliche Begründung auf S.\;\pageref{58952_dat}.%
\pend%
\end{ledgroup}%
\selectlanguage{latin}%
\frenchspacing%
\vspace{8mm}%
\normalsize%
\pstart%
\noindent%
\lbrack%
\href{\Scan{LH_35_11_14_037v}}%
{37\:v\textsuperscript{o}}%
\rbrack\ %
Videamus
an refractionis regula%
\indexS{regula refractionis}%
\indexS{refractio corporum solidorum}
in corporibus solidis%
\indexS{corpus ingrediens liquida}
liquida ingredientibus%
\indexS{liquidum}
explicari possit per corpus%
\indexS{corpus incurrens in globulos}
in globulorum%
\indexS{series globulorum}%
\indexS{globulus}
\edtext{%
series incurrens.%
}{%
\lemma{series}%
\Bfootnote{%
\textit{(1)}~\textlangle ob\textrangle currens. %
\textit{(2)}~incurrens.%
~\textit{L}%
}}
Et quomodo explicari possit
quod experimur,
ut corpora solida%
\indexS{corpus impulsum in aquam}
nimis oblique in aquam impulsa%
\indexS{aqua impulsa a corpore}
non penetrent,%
\indexS{corpus penetrans aquam}
sed resiliant.%
\indexS{corpus resiliens ab aqua}
Videndum
quid fiat
\edtext{%
si excipientia sint magnitudinis inassignabilis,%
\indexS{magnitudo inassignabilis}%
\indexS{corpus magnitudinis inassignabilis}%
}{%
\lemma{si}%
\Bfootnote{%
\textit{(1)}~incurrens sit infinit\textlangle ae magni\textrangle\  %
\textit{(2)}~excipientia sint magnitudinis inassignabilis,%
~\textit{L}%
}}
ubi infinitus eorum numerus pellendus%
\indexS{numerus infinitus}
antequam tot­um incurrens immergatur.
\pend%
\begin{edfigure}[2.25em][0.00em]
\includegraphics%
{gesamttex/edit_VIII-4/images/LH_35_11_14_037_d_037v.pdf}
\edcaption[0.5em]{%
\lbrack\textit{Fig.\:1}\rbrack%
}%
\end{edfigure}%
\newpage%
\pstart%
\noindent%
\centering%
*     *     *
\pend%
\vspace{4.0em}%
\selectlanguage{german}%
\frenchspacing%
\footnotesize%
\pstart%
\noindent%
\textls{Ausführliche Datierungsgründe:}%
\pagelabel{58952_dat}
Die titellose und undatierte Notiz N.\;\ref{58952}, ein kurzer Vermerk über Bewegung in flüssigen Medien, ist auf einem Träger überliefert, der kein Wasserzeichen aufweist. %
Die auf Bl.\;\href{\Scan{LH_35_11_14_037r}}{37\:r\textsuperscript{o}} überlieferten Textfragmente lassen sich jedoch einem sonst nicht bekannten Entwurf von Leibnizens Besprechung des unauthorisierten Neudrucks seiner \textit{Dissertatio de arte combinatoria} zuordnen (\textsc{G.\:W.\;Leibniz}, \textit{Ars combinatoria}, Frankfurt am Main 1690).\cite{01485}
Da diese Besprechung im Februarheft der \textit{Acta eruditorum} 1691 (S.\;63\:f.)\cite{01023}\cite{01486}
anonym erschien, müssen die Textfragmente auf Bl.\;\href{\Scan{LH_35_11_14_037r}}{37\:r\textsuperscript{o}} früher entstanden sein, wahrscheinlich schon vor Ende 1690, spätestens aber Anfang 1691.
Die Notiz N.\;\ref{58952} könnte hingegen auch etwas später verfasst worden sein.
Eine viel spätere Entstehungszeit als Anfang 1691 ist freilich nicht ganz auszuschließen, aber unwahrscheinlich.
\pend%
\pstart%
Wann Leibniz vom unauthorisierten Frankfurter\indexO{Frankfurt am Main} Druck seiner \textit{Dissertatio de arte combinatoria} erfuhr, ist nicht bekannt.
Als er in seinem Brief an Daniel Georg Morhof\indexN{\textls{Morhof} (Morhofius), Daniel Georg 1639\textendash1691} vom August 1690 sein Jugendwerk ausdrücklich erwähnt (\textit{LSB} II,\;2 N.\;82, S.\;332.27--333.2),\cite{01487} verweist er, anders als in Briefen späterer Jahre, mit keinem Wort auf den Neudruck.
Dies lässt vermuten, dass er im August 1690 noch nichts von dem unauthoriserten Druck wusste und die Besprechung für die \textit{Acta eruditorm} folglich nicht davor verfassen konnte.
August 1690 ist daher als Terminus post quem der Datierung der vorliegenden Notiz N.\;\ref{58952} anzusehen.
\pend%
\selectlanguage{latin}%
\frenchspacing%
\normalsize%
\newpage%
